% Part: first-order-logic
% Chapter: completeness
% Section: outline

\documentclass[../../../include/open-logic-section]{subfiles}

\begin{document}

\iftag{FOL}
      {\olfileid{fol}{com}{out}}
      {\olfileid{pl}{com}{out}}

\olsection{প্রমাণের রূপরেখা}

পূর্ণতা উপপাদ্যের প্রমাণ কিছুটা জটিল; প্রথমবার পড়ার সময় পথ হারানো সহজ।
তাই প্রমাণটির একটি রূপরেখা দিই। প্রথম ধাপটি দৃষ্টিভঙ্গির পরিবর্তন; এর ফলে
প্রমাণের একটি পথ দেখা যায়। পূর্ণতাকে যখন “$\Gamma \Entails !A$ হলেই
$\Gamma \Proves !A$” হিসেবে ভাবি, তখন কোনো ধারণা পাওয়াই কঠিন হতে পারে:
কারণ $\Gamma \Proves !A$ দেখাতে আমাদের একটি !!a{derivation} খুঁজতে হয়,
অথচ $\Gamma \Entails !A$ পূর্বধারণাটি এই কাজে কোনোভাবে সাহায্য করে বলে
মনে হয় না। কিছু প্রমাণপদ্ধতির ক্ষেত্রে সরাসরি একটি !!a{derivation}
নির্মাণ করা সম্ভব, কিন্তু আমরা একটু ভিন্ন পথ নেব। প্রয়োজনীয় দৃষ্টিভঙ্গির
পরিবর্তনটি হলো: পূর্ণতাকে এভাবেও বলা যায়—“$\Gamma$ সঙ্গত হলে সেটি
পরিতৃপ্তিযোগ্য।” হয়তো~$\Gamma$-র তথ্য ও তার সঙ্গতির পূর্বধারণা একসঙ্গে
ব্যবহার করে এমন একটি \iftag{FOL}{!!a{structure}}{!!a{valuation}} নির্মাণ
করতে পারি, যা~$\Gamma$-র প্রতিটি
\iftag{FOL}{!!{sentence}}{!!{formula}}-কে পরিতৃপ্ত করে। আমরা তো জানিই
কেমন \iftag{FOL}{!!{structure}}{!!{valuation}} চাই: $\Gamma$ তাকে যেরূপ
বর্ণনা করে, ঠিক সেরূপ!{}

$\Gamma$-তে কেবল \iftag{FOL}{পরমাণু !!{sentence}s}{!!{propositional variable}s}
থাকলে তার একটি মডেল নির্মাণ করা সহজ।\iftag{FOL}{ ধরা যাক,
  পরমাণু !!{sentence}s সবই $\Atom{P}{a_1,\dots,a_n}$ আকারের, যেখানে
  $a_i$-গুলি !!{constant}s।}{} আমাদের শুধু এমন
\iftag{FOL}{%
  একটি !!a{domain}~$\Domain{M}$ এবং~$P$-এর এমন একটি মানায়ন দিতে হবে যাতে
  $\Sat{M}{\Atom{P}{a_1,\ldots,a_n}}$। কাজটি খুব কঠিন নয়:
  $\Domain{M} = \Nat$, $\Assign{\Obj c_i}{M} = i$ রাখি; এবং প্রত্যেক
  $\Atom{P}{a_1,\ldots,a_n} \in \Gamma$-র জন্য $\tuple{k_1,
    \dots, k_n}$ টিউপলটি $\Assign{P}{M}$-তে রাখি, যেখানে $k_i$ হলো
  ধ্রুবক সংকেত~$a_i$-এর সূচক (অর্থাৎ $a_i \ident \Obj c_{k_i}$)।}{%
  এমন একটি !!a{valuation}~$\pAssign{v}$ দিতে হবে যাতে প্রত্যেক $p \in
  \Gamma$-র জন্য $\pSat{v}{p}$। ধরি, $p \in \Gamma$ যদি এবং কেবল যদি
  $\pAssign{v}(p) = \True$।}

এবার ধরা যাক $\Gamma$-তে কোনো !!{formula}~$\lnot !B$ আছে, যেখানে $!B$
পরমাণু। আমাদের আশঙ্কা হতে পারে যে
$\iftag{FOL}{\Struct{M}}{\pAssign{v}}$-র নির্মাণটি $\lnot !B$-কে সত্য
করার সম্ভাবনায় বাধা দেবে। কিন্তু এখানেই~$\Gamma$-র সঙ্গতি কাজে লাগে:
$\lnot !B \in \Gamma$ হলে $!B \notin \Gamma$; নইলে $\Gamma$ অসঙ্গত
হতো। আর $!B \notin \Gamma$ হলে
$\iftag{FOL}{\Struct{M}}{\pAssign{v}}$-র আমাদের নির্মাণ অনুসারে
$\iftag{FOL}{\Sat/{M}}{\pSat/{v}}{!B}$, তাই
$\iftag{FOL}{\Sat{M}}{\pSat{v}}{\lnot !B}$। এত দূর পর্যন্ত ঠিক আছে।

$\Gamma$-তে জটিল, অ-পরমাণু সূত্র থাকলে কী হবে? ধরা যাক, তাতে $!A
\land !B$ আছে। একে সত্য করতে হলে আমাদের এমনভাবে এগোনো উচিত যেন $!A$ ও
$!B$ উভয়ই~$\Gamma$-তে আছে। আর $!A \lor !B \in \Gamma$ হলে অন্তত
একটিকে সত্য করতে হবে; অর্থাৎ এমনভাবে এগোতে হবে যেন তাদের একটি~$\Gamma$-তে
আছে।

এ থেকে নিচের ধারণাটি পাওয়া যায়: আমরা~$\Gamma$-তে আরও !!{formula}s যোগ
করি, যাতে (ক)~ফলে পাওয়া সেটটি সঙ্গত থাকে, (খ)~প্রতিটি সম্ভাব্য পরমাণু
!!{sentence}~$!A$-এর জন্য ফলে পাওয়া সেটটিতে $!A$ অথবা $\lnot !A$
থাকে, এবং (গ)~যখনই সেটটিতে $!A \land !B$ থাকে তখন $!A$ ও $!B$ উভয়ই
থাকে, যখন $!A \lor !B$ থাকে তখন $!A$ বা $!B$-এর অন্তত একটি থাকে,
ইত্যাদি। আমরা এভাবে (সম্ভবত অনন্তকাল) যোগ করতে থাকি। এইভাবে যোগ করা সব
!!{formula}s-এর সেটকে~$\Gamma^*$ বলি। তখন ওপরের নির্মাণটি আমাদের এমন একটি
\iftag{FOL}{!!a{structure}~$\Struct{M}$}{!!a{valuation}~$\pAssign{v}$}
দেবে, যার সম্বন্ধে আরোহের সাহায্যে প্রমাণ করা যাবে যে সেটি~$\Gamma^*$-র
সব বাক্যকে পরিতৃপ্ত করে; এবং $\Gamma \subseteq \Gamma^*$ হওয়ায়
সেটি~$\Gamma$-র সব বাক্যকেও পরিতৃপ্ত করে। দেখা যায়, (ক) ও~(খ)-এর
নিশ্চয়তাই যথেষ্ট। যে বাক্য-সেটের জন্য (খ) সত্য, তাকে \emph{পূর্ণ} বলা
হয়। কাজেই আমাদের কাজ হবে সঙ্গত সেট~$\Gamma$-কে একটি সঙ্গত ও পূর্ণ
সেট~$\Gamma^*$-তে প্রসারিত করা।

\iftag{FOL}{%
এই পরিকল্পনায় একটি জটিলতা আছে: $\lexists[x][!A(x)] \in \Gamma$ হলে
আমরা কোনো !!{constant}~$c$ বেছে নিয়ে এই প্রক্রিয়ায় $!A(c)$ যোগ করতে
চাইব। কিন্তু সব সময় তা করা যাবে কীভাবে জানব? হয়তো আমাদের ভাষায় মাত্র
কয়েকটি !!{constant}s আছে, এবং তাদের প্রত্যেকটির জন্য $\lnot !A(c) \in
\Gamma$। তখন $!A(c)$-ও যোগ করতে পারি না, কারণ তাতে সেটটি অসঙ্গত হবে;
আর $\Struct{M}$-কে $!A(c)$ নাকি $\lnot !A(c)$ সত্য করতে হবে, তাও জানা
থাকবে না। এমনও হতে পারে যে~$\Gamma$ এমন একটি ভাষার কেবল বাক্যগুলি নিয়ে
গঠিত, যেখানে কোনো ধ্রুবক সংকেতই নেই (যেমন, সেটতত্ত্বের ভাষা)।

এই সমস্যার সমাধান হলো শুরুতেই অসীমসংখ্যক ধ্রুবক এবং তাদের সঙ্গে
পরিমাণসূচকগুলিকে উপযুক্তভাবে যুক্ত করে এমন বাক্যগুলি যোগ করা। (অবশ্যই
যাচাই করতে হবে যে এতে অসঙ্গতি আসতে পারে না।)

পরমাণু বাক্যে কেবল !!{constant}s থাকলে আমাদের আদি নির্মাণ ভালো কাজ করে।
কিন্তু ভাষায় !!{function}s-ও থাকতে পারে। সেক্ষেত্রে সবকিছু কার্যকর করতে
এই !!{function}s-কে দেওয়ার জন্য~$\Nat$-এর ওপর ঠিক অপেক্ষকগুলি খুঁজে পাওয়া
কঠিন হতে পারে। তাই আরেকটি কৌশল নিই: $\Obj c_i$-র ব্যাখ্যা হিসেবে $i$
ব্যবহার না করে, !!{constant}s-এর সেটটিকেই সংজ্ঞাক্ষেত্র করি। তাহলে
$\Struct M$ প্রতিটি !!{constant}-কে নিজেকেই দিতে পারে:
$\Assign{\Obj c_i}{M} = \Obj c_i$। আরও এক ধাপ এগোই না কেন:
$\Domain{M}$-কে ভাষাটির সব \emph{পদ}-এর সেট করি!{} এমন করলে
!!{function}s-কে দেওয়ার জন্য অপেক্ষকের একটি স্পষ্ট মানায়ন আছে—এই
অপেক্ষকগুলির যুক্তি ও মান উভয়ই পদ। !!{function}~$\Obj f^n_i$-কে এমন
অপেক্ষকটি দিই, যা নিবেশ হিসেবে $n$-টি পদ $t_1$, \dots,~$t_n$ পেলে মান
হিসেবে $\Obj f^n_i(t_1, \dots, t_n)$ পদটি দেয়।

ধাঁধার শেষ অংশ হলো~$\eq$ নিয়ে কী করা হবে। !!{predicate}~$\eq$-এর
ব্যাখ্যা স্থির: $\Sat{M}{\eq[t][t']}$ যদি এবং কেবল যদি
$\Value{t}{M} = \Value{t'}{M}$। এখন যদি এমনভাবে সবকিছু সাজাই যে
পদ~$t$-এর !!{value} নিজেই $t$, তবে এই !!{structure} $\eq[t][t']$
আকারের কোনো বাক্যকে সত্য করবে না, যদি না $t$ ও $t'$ একই পদ হয়। অবশ্যই
এটি একটি সমস্যা; কারণ !!{function}s-সহ ভাষায় প্রায় প্রতিটি আকর্ষণীয়
তত্ত্বেই এমন $\eq[t][t']$ বাক্য উপপাদ্য হয়, যেখানে $t$ ও $t'$ একই পদ
নয় (যেমন, পাটীগণিতের তত্ত্বে: $\eq[(\Obj 0+ \Obj 0)][\Obj 0]$)। এই
সমস্যা সমাধান করতে~$\Struct M$-এর সংজ্ঞাক্ষেত্র বদলাই: $\Domain{M}$-এর
বস্তু হিসেবে পদ ব্যবহার না করে পদগুলির সেট ব্যবহার করি; প্রতিটি সেটে সেই
সব পদ থাকবে, যেগুলিকে~$\Gamma$-র বাক্যগুলি সমান হতে বাধ্য করে। যেমন,
$\Gamma$ পাটীগণিতের কোনো তত্ত্ব হলে এই সেটগুলির একটিতে থাকবে: $\Obj 0$,
$(\Obj 0 + \Obj 0)$, $(\Obj 0 \times \Obj 0)$, ইত্যাদি। $\Obj 0$-কে
আমরা এই সেটটিই দেব; দেখা যাবে, এই সেটটি এর অন্তর্ভুক্ত সব পদেরও মান—যেমন,
$(\Obj 0 + \Obj 0)$-এরও। তাই সংশোধিত !!{structure}-টিতে
$\eq[(\Obj 0+ \Obj 0)][\Obj 0]$ বাক্যটি সত্য হবে।}{}

তাহলে আমরা যা করব তা হলো এই। প্রথমে !!{complete} সঙ্গত সেটগুলির ধর্ম
অনুসন্ধান করব। বিশেষত, প্রমাণ করব যে একটি !!a{complete} সঙ্গত সেটে
$!A \land !B$ থাকে যদি এবং কেবল যদি $!A$ ও~$!B$ উভয়ই থাকে; $!A \lor
!B$ থাকে যদি এবং কেবল যদি তাদের অন্তত একটি থাকে; ইত্যাদি
(\olref[ccs]{prop:ccs})।\iftag{FOL}{ তারপর বাক্যের “সম্পৃক্ত” সেটের
  সংজ্ঞা দিয়ে তার ধর্ম অনুসন্ধান করব। সম্পৃক্ত সেটে এমন শর্তমূলক বাক্য
  থাকে, যা প্রতিটি পরিমাণিত !!{sentence}-কে তার নিদর্শনগুলির সঙ্গে যুক্ত
  করে (\olref[hen]{defn:henkin-exp})। আমরা দেখাব, যেকোনো সঙ্গত
  সেট~$\Gamma$-কে সব সময় একটি সম্পৃক্ত সেট~$\Gamma'$-তে প্রসারিত করা
  যায় (\olref[hen]{lem:henkin})। কোনো সেট সঙ্গত, সম্পৃক্ত ও !!{complete}
  হলে তার আরও এই ধর্ম থাকে যে সেটিতে
  \iftag{prvEx}{$\lexists[x][!A(x)]$ থাকে যদি এবং কেবল যদি কোনো বদ্ধ
    পদ~$t$-এর জন্য $!A(t)$ থাকে}{}\iftag{defEx,defAll}{}{ এবং
  }\iftag{prvAll}{$\lforall[x][!A(x)]$ থাকে যদি এবং কেবল যদি প্রতিটি
    বদ্ধ পদ~$t$-এর জন্য $!A(t)$ থাকে}{} (\olref[hen]{prop:saturated-instances})।}{}
তারপর \iftag{FOL}{সম্পৃক্ত}{} সঙ্গত সেট~$\iftag{FOL}{\Gamma'}{\Gamma}$
নিয়ে দেখাব যে একে একটি \iftag{FOL}{সম্পৃক্ত, সঙ্গত ও}{সঙ্গত ও}
!!{complete} সেট~$\Gamma^*$-তে প্রসারিত করা যায়
(\olref[lin]{lem:lindenbaum})। এই $\Gamma^*$ দিয়েই আমাদের
\iftag{FOL}{পদ-মডেল~$\Struct M(\Gamma^*)$}{মানায়ন~$\pAssign
v(\Gamma^*)$} সংজ্ঞায়িত করব। \iftag{FOL}{পদ-মডেলের সংজ্ঞাক্ষেত্র হলো
  বদ্ধ পদগুলির সেট, আর তার !!{predicate}s-এর ব্যাখ্যা পরমাণু
  !!{sentence}s দিয়ে}{!!{valuation}-টি !!{propositional
    variable}s দিয়ে} $\Gamma^*$ থেকে নির্ধারিত হয়
(\olref[mod]{defn:termmodel})। \iftag{FOL}{সম্পৃক্ত,}{} পূর্ণ সঙ্গত
সেটের ধর্মগুলি ব্যবহার করে দেখাব যে সত্যিই
$\iftag{FOL}{\Sat{M(\Gamma^*)}}{\pSat{v(\Gamma^*)}}{!A}$ যদি এবং কেবল
যদি $!A \in \Gamma^*$ (\olref[mod]{lem:truth}); ফলে বিশেষত,
$\iftag{FOL}{\Sat{M(\Gamma^*)}}{\pSat{v(\Gamma^*)}}{\Gamma}$।
\iftag{FOL}{শেষে, $\Gamma$-তে~$\eq$ থাকলে পদ-মডেল কীভাবে সংজ্ঞায়িত করতে
  হবে তা বিবেচনা করব (\olref[ide]{defn:term-model-factor}) এবং দেখাব যে
  সেটি~$\Gamma^*$-কে পরিতৃপ্ত করে (\olref[ide]{lem:truth})।}{}

\end{document}
